\section{产品经理底盘：第三方视角、体验和数据}

上一章说明这期要读产品判断，本章先看判断能力的来源。明超平大学打辩论，尤其擅长四辩；他总结辩论最重要的不是说服对手，而是永远具备第三方视角。这个训练后来变成产品方法：用户不是团队内部争论的对手，而是台下评委。产品打开的那一瞬间，用户不知道你的故事、战略和辛苦，他只知道自己能不能立刻理解和使用。

\begin{importantbox}{课堂提示：一秒钟把自己变傻}
明超平引用张小龙式的产品直觉：要把自己变成第一次进入场景的普通用户。任何需要长说明、弹窗和复杂教育才能使用的软件，都没有真正完成体验设计。
\end{importantbox}

\subsection{体验不是数据}

前面讲第三方视角，本节进入 OnePlus 给他的第一课：产品不是指标集合，而是用户在具体情境中的连续感受。这里要回答的问题是，为什么一个看起来被数据证明很好的功能，仍然可能让用户焦虑。手机续航有客观数据，但续航体验不等于电量曲线。用户最焦虑的可能是 95--100\% 的充电等待，或者 0--5\% 的濒死状态；纸面续航漂亮，不代表体验好。这个案例说明，数据能告诉你发生了什么，但不能自动告诉你用户如何感受。

\begin{figure}[H]
\centering
\includegraphics[width=0.94\textwidth]{experience-not-data.png}
\caption{体验不是数据：续航数据和续航体验不划等号。自制概念图，依据 00:26:37--00:33:58 对谈内容整理。}
\end{figure}

\begin{knowledgebox}{读图：体验是数据和心理临界点的组合}
左边 Data 是平均续航、留存、点击等指标；右边 Experience 是焦虑、敏感点和主观感受。读这张图时要注意，体验不是反数据，而是要求你知道哪些数据区间对用户心理最敏感。
\end{knowledgebox}

\begin{knowledgebox}{术语消化：数据、体验、后视镜}
\noindent\begin{tabularx}{\textwidth}{p{0.18\textwidth}p{0.36\textwidth}X}
\toprule
概念 & 含义 & 本期中的判断 \\
\midrule
数据 & 对过去行为的量化记录 & 能验证结果，但不能直接产生方向。 \\
体验 & 用户在具体临界点上的主观感受 & 决定产品是否被理解、信任和继续使用。 \\
后视镜 & 数据像开车后视镜 & 有助于判断过去路径，但不能替你看前方道路。 \\
信仰判断 & 数据短期不支持但方向仍值得押注 & AI 应用窗口里尤其重要。 \\
\bottomrule
\end{tabularx}
\end{knowledgebox}

\subsection{字节的长处和副作用}

从 OnePlus 的体验训练进入字节，问题变成另一套速度和规模。本节看字节给他的第二课：数据驱动和高频迭代。OnePlus 一年做一款手机，字节可能两周一个版本，产品经理必须在数据、实验、增长和组织协同中快速转动。明超平承认字节让他掌握了更大规模的产品方法，但也指出副作用：强数据体系可能磨灭灵光一现的创意，因为某些真正新东西短期数据可能是负向的。

\begin{warningbox}{数据基建不能替代创意判断}
如果只追求短期 A/B 正向，团队会越来越擅长优化已有局部，却不一定能创造新形态。AI 应用创业尤其如此：新用户、新人群和新工作流一开始往往没有稳定数据。
\end{warningbox}

\subsection{智能车竞赛：早期 reward model 训练}

在进入职业产品经理之前，明超平还提到大学后两年基本睡在实验室里做智能汽车竞赛。这个经历很适合放进 AI 应用创业语境里看：小车跑得更快、避障更稳定、通信更顺，都会立刻反馈到比赛成绩上。他把这称为很好的 reward model：代码、传感器、硬件和团队协作的改动，会以非常直接的方式回到结果。

这段经历解释了他后来为什么对“环境”和“反馈”敏感。很多产品经理只在屏幕和数据面板上做判断，而智能车竞赛让他早早接触到真实世界里的闭环：传感器输入、控制输出、bug、速度、比赛结果。YouWare 虽然是软件产品，但它也在寻找类似闭环：用户输入意图，系统生成可交互作品，作品被分享和反馈，下一轮产品再调整。

\begin{importantbox}{实践经验：好的产品反馈要像比赛计时器}
如果一个系统的反馈太慢、太模糊，团队很难形成直觉。智能车竞赛的强反馈让人上瘾，AI 应用也需要类似机制：用户是否发布、是否分享、是否被别人互动、是否继续创作，都比单次点击更接近真实 reward。
\end{importantbox}

\subsection{本章小结}

明超平的产品底盘来自四种训练：辩论训练第三方视角，智能车竞赛训练闭环反馈，OnePlus 训练体验敏感，字节训练数据和迭代。这些能力合起来，解释了他后来为什么强调容器、环境和 per token valuation。

